\documentclass{article}
\usepackage{amsmath, amssymb, amsthm}
\usepackage{hyperref}
\usepackage{geometry}
\geometry{margin=1in}

\title{Erd\H{o}s Problem \#317}
\date{Last updated: 2025-08-31}
\author{Source: erdosproblems.com}

\begin{document}
\maketitle

\noindent\textbf{Status:} OPEN \quad \textbf{No prize}

\noindent\textbf{Tags:} number theory, unit fractions

\medskip

\noindent\textbf{Problem Statement:}

Is there some constant $c&#62;0$ such that for every $n\geq 1$ there exists some $\delta_k\in \{-1,0,1\}$ for $1\leq k\leq n$ with\[0&lt; \left\lvert \sum_{1\leq k\leq n}\frac{\delta_k}{k}\right\rvert &lt; \frac{c}{2^n}?\]Is it true that for sufficiently large $n$, for any $\delta_k\in \{-1,0,1\}$,\[\left\lvert \sum_{1\leq k\leq n}\frac{\delta_k}{k}\right\rvert &gt; \frac{1}{[1,\ldots,n]}\]whenever the left-hand side is not zero?




Inequality is obvious for the second claim, the problem is strict inequality. This fails for small $n$, for example\[\frac{1}{2}-\frac{1}{3}-\frac{1}{4}=-\frac{1}{12}.\]Arguments of Kovac and van Doorn in the comment section prove a weak version of the first question, with an upper bound of\[2^{-n\frac{(\log\log\log n)^{1+o(1)}}{\log n}},\]and van Doorn gives a heuristic that suggests this may be the true order of magnitude.


\bigskip
\noindent\small{Source: \url{https://www.erdosproblems.com/317}}
\end{document}
